\documentclass[11pt]{article}
\usepackage[margin=1in]{geometry}
\usepackage{amsmath,amssymb,amsthm}
\usepackage{graphicx}
\usepackage[hidelinks]{hyperref}
\setlength{\emergencystretch}{2em}

\newtheorem{theorem}{Theorem}
\newtheorem{lemma}{Lemma}
\newtheorem{corollary}{Corollary}
\newtheorem{definition}{Definition}

\newcommand{\N}{\mathbb N}
\newcommand{\cI}{\mathcal I}

\title{Ideal Reduced Products and a Diagonal Obstruction to BCP}
\author{Run \texttt{fa\_banach\_001}, lane 19}
\date{June 17, 2026}

\begin{document}
\maketitle

\section*{Source Problem}

The source paper is Jinghao Huang, Karimbergen Kudaybergenov, and Rui Liu,
\emph{The ball-covering property of von Neumann algebras and noncommutative
symmetric spaces}, arXiv:2311.14257.  After proving that a commutative unital
\(C^*\)-algebra \(C(K)\) has BCP/UBCP exactly when \(K\) has a countable
\(\pi\)-basis, the authors ask for a noncommutative version:

\begin{center}
\includegraphics[width=0.92\textwidth]{figures/open_problem_crop.png}
\end{center}

\section*{Definitions}

An ideal \(\cI\) on \(\N\) is called proper if \(\N\notin\cI\).  A set
\(S\subseteq\N\) is \(\cI\)-positive if \(S\notin\cI\).

\begin{definition}
A proper ideal \(\cI\) on \(\N\) has the countable disjoint-refinement
property if, for every sequence \((S_k)\) of \(\cI\)-positive subsets of
\(\N\), there are pairwise disjoint \(\cI\)-positive sets \(T_k\subseteq S_k\).
\end{definition}

Let \((E_n)\) be a sequence of normed spaces.  Define
\[
        c_\cI(E_n)
        =
        \{x=(x_n)\in \ell_\infty(E_n):
        \{n:\|x_n\|\ge \varepsilon\}\in\cI
        \text{ for every }\varepsilon>0\}.
\]
This is the usual \(c_0\)-sum when \(\cI\) is the ideal of finite subsets of
\(\N\).  It is closed in \(\ell_\infty(E_n)\): if \(x^{(m)}\to x\) in
\(\ell_\infty\) and \(x^{(m)}\in c_\cI(E_n)\), then for fixed
\(\varepsilon>0\) and large \(m\),
\[
        \{n:\|x_n\|\ge\varepsilon\}
        \subseteq
        \{n:\|x^{(m)}_n\|\ge\varepsilon/2\}\in\cI.
\]

\section*{Main Result}

\begin{theorem}
\label{thm:ideal-reduced-products}
Let \(\cI\) be a proper ideal on \(\N\) containing all finite sets and having
the countable disjoint-refinement property.  If \((E_n)\) is any sequence of
nonzero normed spaces, then
\[
        \ell_\infty(E_n)/c_\cI(E_n)
\]
fails BCP.
\end{theorem}

\begin{corollary}
\label{cor:cstar-reduced-products}
Let \(\cI\) be as in Theorem~\ref{thm:ideal-reduced-products}.  If
\((A_n)\) is any sequence of nonzero \(C^*\)-algebras, then the
\(C^*\)-reduced product
\[
        \prod_{n=1}^\infty A_n \Big/ \bigoplus\nolimits_{\cI} A_n
\]
fails BCP.
\end{corollary}

\section*{Idea of the Proof}

The ordinary product quotient \(\ell_\infty(E_n)/c_0(E_n)\) has a
Fréchet-limsup norm.  The previous diagonal proof attacks different proposed
centers along different infinite coordinate sets.  For a general ideal
\(\cI\), the same proof works exactly when the ideal lets us refine countably
many \(\cI\)-positive near-norm sets into pairwise disjoint \(\cI\)-positive
attack sets.

On the attack set assigned to a center, choose unit vectors in the fibers
pointing far enough away from that center coordinate.  The assembled bounded
sequence has reduced-product norm one, while each proposed center sees a
positive set of coordinates on which the distance is at least its own norm.

\section*{Proof}

We use the standard equivalent form of BCP: if a normed space \(Y\) has BCP,
then its unit sphere is covered by countably many open balls
\[
        B(y_j,\|y_j\|).
\]
Thus, to prove failure of BCP, it is enough to show that every countable family
of centers misses some point of the unit sphere.

\begin{lemma}[reduced-product norm]
\label{lem:reduced-norm}
For \(x=(x_n)\in\ell_\infty(E_n)\),
\[
        \|x+c_\cI(E_n)\|
        =
        \inf_{A\in\cI}\ \sup_{n\notin A}\|x_n\|.
\]
\end{lemma}

\begin{proof}
Let \(L\) denote the right-hand side.  If \(y\in c_\cI(E_n)\) and
\(\varepsilon>0\), then \(A_\varepsilon=\{n:\|y_n\|\ge\varepsilon\}\in\cI\).
Hence
\[
        \|x-y\|_\infty
        \ge
        \sup_{n\notin A_\varepsilon}\|x_n-y_n\|
        \ge
        \sup_{n\notin A_\varepsilon}\|x_n\|-\varepsilon
        \ge L-\varepsilon.
\]
Taking the infimum over \(y\) and then letting \(\varepsilon\downarrow0\)
gives the lower bound.

Conversely, for \(A\in\cI\), define \(y_n=x_n\) on \(A\) and \(y_n=0\) off
\(A\).  Since the support of \(y\) belongs to \(\cI\), we have
\(y\in c_\cI(E_n)\), and
\[
        \|x-y\|_\infty=\sup_{n\notin A}\|x_n\|.
\]
Taking the infimum over \(A\in\cI\) gives the reverse inequality.
\end{proof}

\begin{lemma}[one-coordinate escape]
\label{lem:coordinate-escape}
Let \(E\ne\{0\}\) be a normed space and \(y\in E\).  There is
\(u\in S_E\) such that
\[
        \|u-y\|\ge \|y\|.
\]
\end{lemma}

\begin{proof}
If \(y=0\), take any \(u\in S_E\).  If \(y\ne0\), choose \(f\in S_{E^*}\) with
\(f(y)=\|y\|\).  Since \(\|f\|=1\), after changing phase if necessary choose
\(u\in S_E\) with \(\operatorname{Re}f(u)<-1/2\).  Then
\[
        \|u-y\|\ge \operatorname{Re}f(y-u)>\|y\|+\frac12\ge\|y\|.
\]
\end{proof}

\begin{lemma}
\label{lem:fin-has-refinement}
The ideal of finite subsets of \(\N\) has the countable disjoint-refinement
property.
\end{lemma}

\begin{proof}
Let \((S_k)\) be a sequence of infinite subsets of \(\N\).  Inductively, at
stage \(m\), choose distinct fresh points \(t_{k,m}\in S_k\) for
\(1\le k\le m\), avoiding all points chosen in previous stages and earlier in
the same stage.  This is possible because only finitely many points are
forbidden when each choice is made.  Then
\[
        T_k=\{t_{k,m}:m\ge k\}
\]
are pairwise disjoint infinite subsets with \(T_k\subseteq S_k\).
\end{proof}

\begin{proof}[Proof of Theorem~\ref{thm:ideal-reduced-products}]
Let
\[
        \xi_i=a^i+c_\cI(E_n),\qquad i\ge1,
\]
be an arbitrary countable family of centers in the reduced product, where
\(a^i=(a^i_n)\in\ell_\infty(E_n)\).  Write \(r_i=\|\xi_i\|\).

Choose a map \(\pi:\N\to\N\) such that every fiber \(\pi^{-1}(\{i\})\) is
infinite.  For each \(k\), set \(i=\pi(k)\) and define
\[
        \delta_k=
        \begin{cases}
        \min\{r_i/2,1/k\},& r_i>0,\\
        1,& r_i=0,
        \end{cases}
        \qquad
        S_k=\{n:\|a^i_n\|>r_i-\delta_k\}.
\]
The set \(S_k\) is \(\cI\)-positive.  Indeed, if \(S_k\in\cI\), then
Lemma~\ref{lem:reduced-norm} would give
\[
        r_i\le \sup_{n\notin S_k}\|a^i_n\|\le r_i-\delta_k
\]
when \(r_i>0\), a contradiction.  If \(r_i=0\), then \(S_k=\N\), which is
\(\cI\)-positive because \(\cI\) is proper.

By the countable disjoint-refinement property, choose pairwise disjoint
\(\cI\)-positive sets \(T_k\subseteq S_k\).  Define \(b=(b_n)\) as follows.
If \(n\in T_k\), use Lemma~\ref{lem:coordinate-escape} in the nonzero space
\(E_n\) to choose \(b_n\in S_{E_n}\) satisfying
\[
        \|b_n-a^{\pi(k)}_n\|\ge \|a^{\pi(k)}_n\|.
\]
This is unambiguous because the sets \(T_k\) are pairwise disjoint.  For
\(n\notin\bigcup_k T_k\), put \(b_n=0\).

Then \(\|b\|_{\ell_\infty}\le1\), and
\(\|b+c_\cI(E_n)\|=1\): the norm is at most one, while each \(T_k\) is
\(\cI\)-positive and \(\|b_n\|=1\) on \(T_k\).

Fix \(i\).  If \(r_i=0\), then
\[
        \|b+c_\cI(E_n)-\xi_i\|\ge0=r_i
\]
is automatic.  Assume \(r_i>0\).  There are infinitely many \(k\) with
\(\pi(k)=i\).  For such \(k\), on the \(\cI\)-positive set \(T_k\) we have
\[
        \|b_n-a^i_n\|\ge \|a^i_n\|>r_i-\delta_k.
\]
Lemma~\ref{lem:reduced-norm} applied to \(b-a^i\) gives
\[
        \|b+c_\cI(E_n)-\xi_i\|\ge r_i-\delta_k.
\]
Letting \(k\to\infty\) along the infinite set \(\{k:\pi(k)=i\}\), we obtain
\[
        \|b+c_\cI(E_n)-\xi_i\|\ge r_i=\|\xi_i\|.
\]
Thus the unit vector \(b+c_\cI(E_n)\) lies in none of the open balls
\(B(\xi_i,\|\xi_i\|)\).  The countable family of centers was arbitrary, so the
reduced product fails BCP.
\end{proof}

\begin{proof}[Proof of Corollary~\ref{cor:cstar-reduced-products}]
For \(C^*\)-algebras \(A_n\), the space
\[
        \bigoplus\nolimits_{\cI} A_n
        =
        c_\cI(A_n)
\]
is a closed two-sided \(C^*\)-ideal of \(\prod_n A_n\).  Apply
Theorem~\ref{thm:ideal-reduced-products} to the underlying normed spaces
\(E_n=A_n\).
\end{proof}

\section*{Scope}

For the finite ideal, Lemma~\ref{lem:fin-has-refinement} recovers the product
quotient obstruction
\(\ell_\infty(E_n)/c_0(E_n)\).  The present formulation isolates the exact
combinatorial feature used by the diagonal argument: different centers must be
attackable on disjoint positive index sets.

The theorem does not cover ultraproducts over a free ultrafilter.  In that
case the dual maximal ideal has no two disjoint positive sets.  Thus the
ultraproduct case, if negative, requires a genuinely different simultaneous
escape argument rather than this center-by-center diagonal splitting.

\section*{Verification and Search Bounds}

The proof is elementary and uses only the reduced-product norm formula, the
Hahn--Banach one-coordinate escape lemma, and the stated disjoint-refinement
property.  The run indexes were searched for arXiv:2311.14257 together with
``reduced product'', ``ideal'', ``disjoint refinement'', and ``ultraproduct'';
no existing packet for this ideal-reduced-product BCP obstruction was found.

\section*{Human Review Focus}

Review Lemma~\ref{lem:reduced-norm} and the positivity of the sets
\[
        S_k=\{n:\|a^{\pi(k)}_n\|>r_{\pi(k)}-\delta_k\}.
\]
Once those are correct, the rest of the argument is the same diagonal escape
mechanism as in the ordinary product quotient.

\begin{thebibliography}{9}

\bibitem{HKL}
Jinghao Huang, Karimbergen Kudaybergenov, and Rui Liu,
\emph{The ball-covering property of von Neumann algebras and noncommutative
symmetric spaces},
arXiv:2311.14257, 2023.

\end{thebibliography}

\end{document}
